\section{Brauer's Theorem and Modular Systems}

\begin{definition}[\texorpdfstring{\(p\)-regular}{p-regular} and \texorpdfstring{\(p\)-singular}{p-singular} elements]
    Let \(G\) be a finite group and let \(p\) be a prime number.
    An element \(g\in G\) is called \textbf{\(p\)-regular} if
    \[
        p\nmid \operatorname{ord}(g).
    \]
    It is called \textbf{\(p\)-singular} in the sense of these notes if
    \[
        \operatorname{ord}(g)=p^e
    \]
    for some integer \(e\geq 0\). Equivalently, \(g\) is a \(p\)-element.
\end{definition}

\begin{lemma}[\(p\)-regular and \(p\)-singular decomposition]
    Let \(G\) be a finite group. For every \(x\in G\), there exist unique elements
    \(s,t\in G\) such that
    \[
        x=st,
        \qquad
        st=ts,
    \]
    where \(s\) is \(p\)-regular and \(t\) is \(p\)-singular.
\end{lemma}

\begin{proof}
    Let
    \[
        \operatorname{ord}(x)=p^a m,
        \qquad
        \gcd(p,m)=1.
    \]
    Since \(\gcd(p^a,m)=1\), choose integers \(r,q\in \mathbb Z\) such that
    \[
        rm+qp^a=1.
    \]
    Define
    \[
        t=x^{rm},
        \qquad
        s=x^{qp^a}.
    \]
    Then \(s,t\in \langle x\rangle\), so \(st=ts\), and
    \[
        st=x^{rm+qp^a}=x.
    \]
    Moreover, the order of \(t\) divides \(p^a\), so \(t\) is \(p\)-singular, while the
    order of \(s\) divides \(m\), so \(s\) is \(p\)-regular.

    For uniqueness, suppose
    \[
        x=s_1t_1=s_2t_2
    \]
    are two such decompositions, with \(s_i\) \(p\)-regular, \(t_i\) \(p\)-singular, and
    \(s_it_i=t_is_i\). Since \(s_i\) and \(t_i\) have coprime orders, both \(s_i\) and
    \(t_i\) are powers of \(x\). Thus they lie in the cyclic group \(\langle x\rangle\).
    Inside the finite cyclic group \(\langle x\rangle\), the \(p\)-part and the \(p'\)-part
    of an element are unique. Hence
    \[
        s_1=s_2,
        \qquad
        t_1=t_2.
    \]
\end{proof}

\begin{lemma}[Commutator subspace of a group algebra]
    Let \(G\) be a finite group and let \(k\) be a field. Put
    \[
        S(kG)=\operatorname{span}_k\{\,ab-ba\mid a,b\in kG\,\}.
    \]
    Then
    \[
        S(kG)
        =
        \left\{
            \sum_{g\in G}\lambda_g g
            \ \middle|\
            \sum_{g\in C}\lambda_g=0
            \text{ for every conjugacy class } C\subseteq G
        \right\}.
    \]
\end{lemma}

\begin{proof}
    First observe that if \(g,h\in G\), then
    \[
        hgh^{-1}-g
        =
        h(gh^{-1})-(gh^{-1})h
        \in S(kG).
    \]
    Hence differences of conjugate group elements lie in \(S(kG)\).

    Let
    \[
        \alpha=\sum_{g\in G}\lambda_g g\in S(kG).
    \]
    Since each commutator \(ab-ba\) has total coefficient \(0\) on every conjugacy
    class, it follows that
    \[
        \sum_{g\in C}\lambda_g=0
    \]
    for every conjugacy class \(C\).

    Conversely, suppose
    \[
        \alpha=\sum_{g\in G}\lambda_g g
    \]
    satisfies
    \[
        \sum_{g\in C}\lambda_g=0
    \]
    for every conjugacy class \(C\). For each conjugacy class \(C\), choose a
    representative \(c_C\in C\). Then
    \[
        \sum_{g\in C}\lambda_g g
        =
        \sum_{g\in C}\lambda_g(g-c_C),
    \]
    because \(\sum_{g\in C}\lambda_g=0\). Since each \(g-c_C\) is a difference of
    conjugate elements, it lies in \(S(kG)\). Therefore
    \[
        \sum_{g\in C}\lambda_g g\in S(kG)
    \]
    for every conjugacy class \(C\), and hence
    \[
        \alpha\in S(kG).
    \]
    This proves the claim.
\end{proof}

\begin{theorem}[Brauer]
    Let \(G\) be a finite group, and let \(k\) be a splitting field for \(G\) with
    \[
        \operatorname{char} k=p>0.
    \]
    Then the number of non-isomorphic simple \(kG\)-modules is equal to the number
    of conjugacy classes of \(p\)-regular elements of \(G\).
\end{theorem}

\begin{proof}
    Put
    \[
        A=kG.
    \]
    As before, define
    \[
        S(A)=\operatorname{span}_k\{\,ab-ba\mid a,b\in A\,\},
    \]
    and
    \[
        T(A)=\{\,a\in A\mid a^{p^m}\in S(A)
        \text{ for some }m\geq 0\,\}.
    \]
    By the previous proposition,
    \[
        \#\{\,\text{non-isomorphic simple }kG\text{-modules}\,\}
        =
        \dim_k kG-\dim_k T(A).
    \]
    Thus it remains to show that
    \[
        \dim_k kG-\dim_k T(A)
    \]
    is the number of conjugacy classes of \(p\)-regular elements.

    Let
    \[
        \{x_i\}_{i\in I}
    \]
    be a set of representatives of the conjugacy classes of \(p\)-regular elements of
    \(G\). We claim that
    \[
        \{\,x_i+T(A)\mid i\in I\,\}
    \]
    is a \(k\)-basis of \(kG/T(A)\).

    First we prove that these elements span \(kG/T(A)\). Let \(x\in G\). By the
    \(p\)-regular and \(p\)-singular decomposition, write
    \[
        x=st,
        \qquad
        st=ts,
    \]
    where \(s\) is \(p\)-regular and \(t\) is \(p\)-singular. Since \(t\) has \(p\)-power
    order, for sufficiently large \(n\) we have
    \[
        t^{p^n}=1.
    \]
    Hence, in characteristic \(p\),
    \[
        (x-s)^{p^n}
        =
        \bigl(s(t-1)\bigr)^{p^n}
        =
        s^{p^n}(t-1)^{p^n}
        =
        s^{p^n}(t^{p^n}-1)
        =
        0.
    \]
    Therefore
    \[
        x-s\in T(A).
    \]
    Since \(s\) is \(p\)-regular, it is conjugate to some \(x_i\). Hence
    \[
        s-x_i\in S(A)\subseteq T(A).
    \]
    Therefore
    \[
        x+T(A)=x_i+T(A).
    \]
    Since the group elements form a \(k\)-basis of \(kG\), this proves that the classes
    \(x_i+T(A)\) span \(kG/T(A)\).

    It remains to prove linear independence. Suppose
    \[
        \sum_{i\in I}\lambda_i x_i\in T(A).
    \]
    Then, by the definition of \(T(A)\), there exists \(m\geq 0\) such that
    \[
        \left(\sum_{i\in I}\lambda_i x_i\right)^{p^m}\in S(A).
    \]
    Enlarging \(m\) if necessary, we may assume that
    \[
        x_i^{p^m}=x_i
        \qquad
        \text{for all }i\in I.
    \]
    Indeed, the orders of the \(x_i\) are prime to \(p\), so one can choose \(m\)
    sufficiently large such that
    \[
        p^m\equiv 1\pmod{\operatorname{ord}(x_i)}
    \]
    for all \(i\).

    Using the congruence
    \[
        \left(\sum_i \lambda_i x_i\right)^{p^m}
        \equiv
        \sum_i \lambda_i^{p^m}x_i^{p^m}
        \pmod{S(A)},
    \]
    we get
    \[
        \sum_{i\in I}\lambda_i^{p^m}x_i\in S(A).
    \]
    By the preceding lemma, the sum of coefficients on each conjugacy class must be
    zero. Since the \(x_i\) lie in distinct conjugacy classes, this gives
    \[
        \lambda_i^{p^m}=0
        \qquad
        \text{for all }i.
    \]
    Hence
    \[
        \lambda_i=0
        \qquad
        \text{for all }i.
    \]
    Therefore the classes \(x_i+T(A)\) are linearly independent.

    Consequently,
    \[
        \dim_k kG/T(A)
        =
        |I|,
    \]
    which is precisely the number of conjugacy classes of \(p\)-regular elements of
    \(G\). Hence
    \[
        \#\{\,\text{non-isomorphic simple }kG\text{-modules}\,\}
        =
        |I|.
    \]
\end{proof}

